\begin{helper}
Let $G$ be a $\Delta$-regular graph and let $I$ be sampled from
\noderef{RandomIndependentSetSampling} with parameter $\gamma$.  Fix
$r$ and $X\subseteq N_G(r)$, and assume that no two distinct vertices of
$X$ have a common neighbor outside $N_G[r]$.  If $u,v\in X$ are distinct
and non-adjacent, then, writing
$c_{uv}=|N_G(u)\cap N_G(v)|$,
\[
\Pr[u\in I\hbox{ and }v\in I]
=
\frac{2}{\Delta^2}
\int_{0\le x\le y\le\gamma}
\left(1-\frac{x}{\Delta}\right)^{\Delta-c_{uv}}
\left(1-\frac{y}{\Delta}\right)^\Delta\,dy\,dx .
\]
\end{helper}
\begin{proof}
By \noderef{RandomIndependentSetSampling}, the random output is obtained by
first activating each vertex independently with probability $\gamma/\Delta$,
then assigning independent uniform priorities, and then retaining exactly the
activated vertices whose activated neighbors all have smaller priority.  Thus
the event $\{u,v\subseteq I\}$ is the event that both $u$ and $v$ are
activated and neither is defeated by an activated neighbor of larger priority.

Condition on the activations of $u$ and $v$.  The two activation probabilities
contribute the factor $(\gamma/\Delta)^2$.  Write
$x=\gamma(1-\pi(u))$ and $y=\gamma(1-\pi(v))$.  Since the priorities are
uniform and independent, the pair $(x,y)$ is uniformly distributed on
$[0,\gamma]^2$; after the activation factor is included this gives density
$dx\,dy/\Delta^2$ on the priority square.  The diagonal $x=y$ has measure zero,
so it may be assigned to either chamber without changing the probability.

On the chamber $x\le y$, equivalently $\pi(u)\ge\pi(v)$, a neighbor adjacent
only to $u$ defeats $u$ precisely when it is activated and has priority larger
than $\pi(u)$, which has conditional probability $x/\Delta$.  Such a neighbor
therefore contributes the survival factor $1-x/\Delta$.  A neighbor adjacent
to $v$ defeats $v$ if it is activated with priority larger than $\pi(v)$,
which has conditional probability $y/\Delta$, so it contributes
$1-y/\Delta$.  A common neighbor of $u$ and $v$ defeats at least one of the two
vertices on this chamber exactly when it is activated with priority larger
than $\pi(v)$; it therefore also contributes $1-y/\Delta$.

The hypotheses give the required counts.  Since $G$ is $\Delta$-regular,
each of $u$ and $v$ has exactly $\Delta$ neighbors.  There are
$c_{uv}$ common neighbors, hence $\Delta-c_{uv}$ neighbors of $u$ that are not
neighbors of $v$, and $\Delta-c_{uv}$ neighbors of $v$ that are not neighbors
of $u$.  The clean-neighborhood assumption ensures that every common neighbor
of the independent pair $u,v\in X$ lies in the closed neighborhood of $r$, so
there is no additional outside common-neighbor correlation beyond this actual
intersection count.  Multiplying the independent survival factors on
$x\le y$ gives
\[
\left(1-\frac{x}{\Delta}\right)^{\Delta-c_{uv}}
\left(1-\frac{y}{\Delta}\right)^{c_{uv}}
\left(1-\frac{y}{\Delta}\right)^{\Delta-c_{uv}}
=
\left(1-\frac{x}{\Delta}\right)^{\Delta-c_{uv}}
\left(1-\frac{y}{\Delta}\right)^\Delta .
\]

The chamber $y\le x$ is identical after interchanging $u$ and $v$ and then
renaming the two variables.  Therefore the two ordered chambers contribute two
equal copies of the displayed integral over $0\le x\le y\le\gamma$.  This is
exactly the asserted formula for the probability that the fixed independent
pair $u,v$ survives.
\end{proof}
